\begin{proposition}[半直积的 Lie 括号]
对 $SO(1,d-1)\ltimes\RR^{1,d-1}$ 的群乘法约定，\[[(\omega_1,a_1),(\omega_2,a_2)]=([\omega_1,\omega_2],\omega_1a_2-\omega_2a_1).\]
\end{proposition}

\begin{proof}
使用群乘法 $(\Lambda,a)(\Gamma,b)=(\Lambda\Gamma,a+\Lambda b)$。
取经过单位元的曲线
\[
g_1(t)=(e^{t\omega_1},ta_1),\qquad
g_2(s)=(e^{s\omega_2},sa_2).
\]
它们的切向量为 $(\omega_1,a_1)$、$(\omega_2,a_2)$；计算 Lie 括号只需这些一阶切向量。
第一步乘积为
\[
g_1g_2=(e^{t\omega_1}e^{s\omega_2},ta_1+s e^{t\omega_1}a_2),
\]
而逆元是
\[
g_1^{-1}=(e^{-t\omega_1},-t e^{-t\omega_1}a_1),\qquad
g_2^{-1}=(e^{-s\omega_2},-s e^{-s\omega_2}a_2).
\]
令 $C=e^{t\omega_1}e^{s\omega_2}e^{-t\omega_1}$，则
\[
g_1g_2g_1^{-1}=(C,ta_1+s e^{t\omega_1}a_2-tCa_1).
\]
必须保留这里的 $sa_2$，它在最后一步才与逆元中的平移抵消。
矩阵展开给出
\[
C=I+s\omega_2+ts[\omega_1,\omega_2]+O(s^2,t^2s).
\]
所以 $ts$ 共轭项已经在这一阶段出现，不能先删去再在最后凭空补回。
右乘 $g_2^{-1}$ 后，记 $L=Ce^{-s\omega_2}$，得到
\[
g_1g_2g_1^{-1}g_2^{-1}
=(L,ta_1+s e^{t\omega_1}a_2-tCa_1-sLa_2).
\]
完整平移部分展开为
\[
\begin{aligned}
ta_1+s(a_2+t\omega_1a_2)
-t(a_1+s\omega_2a_1)-sa_2
=ts(\omega_1a_2-\omega_2a_1),
\end{aligned}
\]
余项为 $O(t^2s,ts^2)$。
Lorentz 部分为 $L=I+ts[\omega_1,\omega_2]+O(t^2s,ts^2)$。
故整个交换子满足
\[
g_1g_2g_1^{-1}g_2^{-1}
=\bigl(I+ts[\omega_1,\omega_2],ts(\omega_1a_2-\omega_2a_1)\bigr)
+O(t^2s,ts^2).
\]
与 Lie 括号的混合二阶定义比较，得到
\begin{equation}\label{eq:ms-poincare-unified-bracket}
\boxed{[(\omega_1,a_1),(\omega_2,a_2)]
=([\omega_1,\omega_2],\omega_1a_2-\omega_2a_1).}
\end{equation}
\end{proof}

\begin{proposition}[矢量 Lorentz 生成元的代数]
令 $(J^{\mu\nu})^\alpha{}_{\beta}=\ii(\eta^{\mu\alpha}\delta^\nu{}_{\beta}-\eta^{\nu\alpha}\delta^\mu{}_{\beta})$，则\[\begin{aligned}{}[J^{\mu\nu},J^{\rho\sigma}]={}&-\ii(\eta^{\mu\rho}J^{\nu\sigma}-\eta^{\mu\sigma}J^{\nu\rho}\\&-\eta^{\nu\rho}J^{\mu\sigma}+\eta^{\nu\sigma}J^{\mu\rho}).\end{aligned}\]
\end{proposition}

\begin{proof}
先算矩阵乘积 $J^{\mu\nu}J^{\rho\sigma}$，再交换指标对 $(\mu,\nu)\leftrightarrow(\rho,\sigma)$ 得到 $J^{\rho\sigma}J^{\mu\nu}$，相减读出结构常数。计算虽然直接，但需要仔细跟踪四个 $\eta$-$\delta$ 因子的缩并。

写出矩阵乘积。对矩阵乘积 $(J^{\mu\nu}J^{\rho\sigma})^\alpha{}_\beta$，先对中间指标 $\gamma$ 求和：
\begin{equation*}
(J^{\mu\nu}J^{\rho\sigma})^\alpha{}_\beta=(J^{\mu\nu})^\alpha{}_\gamma(J^{\rho\sigma})^\gamma{}_\beta.
\end{equation*}
代入定义 $(J^{\mu\nu})^\alpha{}_\gamma=i(\eta^{\mu\alpha}\delta^\nu{}_\gamma-\eta^{\nu\alpha}\delta^\mu{}_\gamma)$ 和 $(J^{\rho\sigma})^\gamma{}_\beta=i(\eta^{\rho\gamma}\delta^\sigma{}_\beta-\eta^{\sigma\gamma}\delta^\rho{}_\beta)$：
\begin{equation*}
(J^{\mu\nu}J^{\rho\sigma})^\alpha{}_\beta=i^2\big(\eta^{\mu\alpha}\delta^\nu{}_\gamma-\eta^{\nu\alpha}\delta^\mu{}_\gamma\big)\big(\eta^{\rho\gamma}\delta^\sigma{}_\beta-\eta^{\sigma\gamma}\delta^\rho{}_\beta\big).
\end{equation*}

展开四项并对 $\gamma$ 求和。利用 $\delta^\nu{}_\gamma\eta^{\rho\gamma}=\eta^{\rho\nu}$ 等：
\begin{equation*}
(J^{\mu\nu}J^{\rho\sigma})^\alpha{}_\beta
=-\eta^{\mu\alpha}\eta^{\rho\nu}\delta^\sigma{}_\beta
+\eta^{\mu\alpha}\eta^{\sigma\nu}\delta^\rho{}_\beta
+\eta^{\nu\alpha}\eta^{\rho\mu}\delta^\sigma{}_\beta
-\eta^{\nu\alpha}\eta^{\sigma\mu}\delta^\rho{}_\beta.
\end{equation*}
四项的物理含义：第一项和第三项含 $\delta^\sigma{}_\beta$（输出指标是 $\sigma$），第二项和第四项含 $\delta^\rho{}_\beta$（输出指标是 $\rho$）。

交换 $(\mu,\nu)\leftrightarrow(\rho,\sigma)$ 得 $J^{\rho\sigma}J^{\mu\nu}$。
\begin{equation*}
(J^{\rho\sigma}J^{\mu\nu})^\alpha{}_\beta
=-\eta^{\rho\alpha}\eta^{\mu\sigma}\delta^\nu{}_\beta
+\eta^{\rho\alpha}\eta^{\nu\sigma}\delta^\mu{}_\beta
+\eta^{\sigma\alpha}\eta^{\mu\rho}\delta^\nu{}_\beta
-\eta^{\sigma\alpha}\eta^{\nu\rho}\delta^\mu{}_\beta.
\end{equation*}

相减并整理。把含 $\delta^\sigma{}_\beta$ 的项合并：
\begin{equation*}
(-\eta^{\mu\alpha}\eta^{\rho\nu}+\eta^{\nu\alpha}\eta^{\rho\mu})\delta^\sigma{}_\beta.
\end{equation*}
利用 $\eta$ 对称性改写为 $\eta^{\rho\mu}$ 的系数：第一项 $-\eta^{\mu\alpha}\eta^{\rho\nu}=-\eta^{\rho\nu}\eta^{\mu\alpha}$，第二项 $\eta^{\nu\alpha}\eta^{\rho\mu}=\eta^{\rho\mu}\eta^{\nu\alpha}$。这恰好可以拆成
\begin{equation*}
\eta^{\mu\rho}\big(\eta^{\nu\alpha}\delta^\sigma{}_\beta\big)-\eta^{\mu\sigma}\big(\eta^{\nu\alpha}\delta^\rho{}_\beta\big)-\eta^{\nu\rho}\big(\eta^{\mu\alpha}\delta^\sigma{}_\beta\big)+\eta^{\nu\sigma}\big(\eta^{\mu\alpha}\delta^\rho{}_\beta\big).
\end{equation*}
重组时须保留 $1/i=-i$ 的负号。每一项的括号内正是某个 $J$ 的分量。例如 $\eta^{\nu\alpha}\delta^\sigma{}_\beta-\eta^{\sigma\alpha}\delta^\nu{}_\beta$ 是 $(J^{\nu\sigma})^\alpha{}_\beta/i$。因此
\begin{equation*}
[J^{\mu\nu},J^{\rho\sigma}]^\alpha{}_\beta
=-i\Big(\eta^{\mu\rho}(J^{\nu\sigma})^\alpha{}_\beta-\eta^{\mu\sigma}(J^{\nu\rho})^\alpha{}_\beta-\eta^{\nu\rho}(J^{\mu\sigma})^\alpha{}_\beta+\eta^{\nu\sigma}(J^{\mu\rho})^\alpha{}_\beta\Big).
\end{equation*}
\end{proof}

\begin{proposition}[轨道 Lorentz 生成元的缩并]
若 $\omega_{\mu\nu}=-\omega_{\nu\mu}$、$L^{\mu\nu}=\ii(x^\mu\partial^\nu-x^\nu\partial^\mu)$，则\[\omega^\mu{}_\nu x^\nu\partial_\mu=\tfrac{\ii}{2}\omega_{\mu\nu}L^{\mu\nu}.\]在紧支光滑测试函数域上 $L^{\mu\nu}$ 关于 $L^2(\dd^dx)$ 形式对称。
\end{proposition}

\begin{proof}
轨道部分 $\omega^\mu{}_\nu x^\nu\partial_\mu$ 中，$\omega_{\mu\nu}$ 是反对称参数（$\omega_{\mu\nu}=-\omega_{\nu\mu}$）。

目标是把轨道部分写成 $+\frac{i}{2}\omega_{\mu\nu}L^{\mu\nu}$ 的形式，从而与内禀部分 $\frac{i}{2}\omega_{\mu\nu}\Sigma^{\mu\nu}$ 合并为统一的生成元 $M^{\mu\nu}=L^{\mu\nu}+\Sigma^{\mu\nu}$。

关键工具是反对称参数与对称张量缩并为零。

降指标。轨道部分中 $\omega^\mu{}_\nu$ 带一个上指标和一个下指标。用 Minkowski 度规降上指标：$\omega^\mu{}_\nu=\eta^{\mu\alpha}\omega_{\alpha\nu}$。代入：
\begin{equation*}
\omega^\mu{}_\nu x^\nu\partial_\mu=\eta^{\mu\alpha}\omega_{\alpha\nu}x^\nu\partial_\mu=\omega_{\alpha\nu}x^\nu\partial^\alpha,
\end{equation*}
最后一步把 $\eta^{\mu\alpha}$ 作用到 $\partial_\mu$ 上得到 $\partial^\alpha\equiv\eta^{\mu\alpha}\partial_\mu$。现在 $\omega_{\alpha\nu}$ 是全下指标，反对称性 $\omega_{\alpha\nu}=-\omega_{\nu\alpha}$ 显式可见。

反对称投影。关键观察：$\omega_{\alpha\nu}$ 反对称，而 $x^\nu\partial^\alpha$ 可以分解为对称和反对称两部分：
\begin{equation*}
x^\nu\partial^\alpha=\frac{1}{2}(x^\nu\partial^\alpha+x^\alpha\partial^\nu)+\frac{1}{2}(x^\nu\partial^\alpha-x^\alpha\partial^\nu).
\end{equation*}
对称部分与反对称 $\omega_{\alpha\nu}$ 缩并为零。验证：
\begin{equation*}
\omega_{\alpha\nu}\cdot\frac{1}{2}(x^\nu\partial^\alpha+x^\alpha\partial^\nu)=\frac{1}{2}(\omega_{\alpha\nu}x^\nu\partial^\alpha+\omega_{\alpha\nu}x^\alpha\partial^\nu).
\end{equation*}
第二项交换哑指标 $\alpha\leftrightarrow\nu$：$\omega_{\nu\alpha}x^\nu\partial^\alpha=-\omega_{\alpha\nu}x^\nu\partial^\alpha$（利用 $\omega_{\nu\alpha}=-\omega_{\alpha\nu}$）。因此两项相加为零。只有反对称部分贡献：
\begin{equation*}
\omega_{\alpha\nu}x^\nu\partial^\alpha=\omega_{\alpha\nu}\cdot\frac{1}{2}(x^\nu\partial^\alpha-x^\alpha\partial^\nu).
\end{equation*}

提出因子并定义轨道角动量。把哑指标 $(\alpha,\nu)$ 换回 $(\mu,\nu)$：
\begin{equation*}
\omega^\mu{}_\nu x^\nu\partial_\mu=\frac{1}{2}\omega_{\mu\nu}(x^\nu\partial^\mu-x^\mu\partial^\nu)=-\frac{1}{2}\omega_{\mu\nu}(x^\mu\partial^\nu-x^\nu\partial^\mu).
\end{equation*}
为了与量子力学中 Hermite 生成元的约定一致（无穷小算符 $U=I-\frac{i}{2}\omega_{\mu\nu}M^{\mu\nu}$ 中 $M^{\mu\nu}$ 应为 Hermite），引入因子 $i$：
\begin{equation*}
-\frac{1}{2}\omega_{\mu\nu}(x^\mu\partial^\nu-x^\nu\partial^\mu)=+\frac{i}{2}\omega_{\mu\nu}\cdot i(x^\mu\partial^\nu-x^\nu\partial^\mu).
\end{equation*}
定义轨道角动量算符
\begin{equation}\label{eq:ms-poincare-L-def}
L^{\mu\nu}=i(x^\mu\partial^\nu-x^\nu\partial^\mu).
\end{equation}
于是轨道部分 = $+\frac{i}{2}\omega_{\mu\nu}L^{\mu\nu}$，内禀部分 = $\frac{i}{2}\omega_{\mu\nu}\Sigma^{\mu\nu}$，两者形式统一。

在时空上的 $L^2(\dd^dx)$ 中，先取紧支光滑测试函数。
积分分部给 $\partial^{\nu\dagger}=-\partial^\nu$，而
$[\partial^\nu,x^\mu]=\eta^{\nu\mu}$，故形式伴随满足
\begin{equation*}
(L^{\mu\nu})^\dagger
=\ii(\partial^\nu x^\mu-\partial^\mu x^\nu)
=\ii(x^\mu\partial^\nu-x^\nu\partial^\mu)=L^{\mu\nu}.
\end{equation*}
两个度规项因对称性相消。这证明该测试域上的形式对称性；
具体物理态空间中生成元的自伴域仍需单独确定，不能由此把形式伴随
和自伴性等同起来。
\end{proof}

\begin{proposition}[Clifford 关系给出的旋量生成元]
若 $\{\gamma^\mu,\gamma^\nu\}=2\eta^{\mu\nu}$，则 $\Sigma^{\mu\nu}=\frac{\ii}{4}[\gamma^\mu,\gamma^\nu]$ 满足\[[\Sigma^{\mu\nu},\gamma^\rho]=\ii(\eta^{\nu\rho}\gamma^\mu-\eta^{\mu\rho}\gamma^\nu),\]与 $D^{-1}\gamma^\rho D=\Lambda^\rho{}_\sigma\gamma^\sigma$ 的无穷小形式一致。
\end{proposition}

\begin{proof}
沿用 $D(I+\omega)=I-\frac{\ii}{2}\omega_{\mu\nu}\Sigma^{\mu\nu}$。
协变条件的顺序为 $D^{-1}\gamma^\rho D=\Lambda^\rho{}_\sigma\gamma^\sigma$，
左端展开成
\[
D^{-1}\gamma^\rho D
=\gamma^\rho+\frac{\ii}{2}\omega_{\mu\nu}
[\Sigma^{\mu\nu},\gamma^\rho]+O(\omega^2).
\]
右端利用 $\omega_{\mu\nu}=-\omega_{\nu\mu}$ 写成
\[
\gamma^\rho+\omega^\rho{}_\sigma\gamma^\sigma
=\gamma^\rho+\frac12\omega_{\mu\nu}
(\eta^{\mu\rho}\gamma^\nu-\eta^{\nu\rho}\gamma^\mu).
\]
这里只能在与反对称 $\omega$ 缩并以后把两项合并；
不缩并的指标恒等式仍必须保留 $1/2$。
比较系数得到
\begin{equation}\label{eq:ms-poincare-sigma-gamma-comm}
[\Sigma^{\mu\nu},\gamma^\rho]
=\ii(\eta^{\nu\rho}\gamma^\mu-\eta^{\mu\rho}\gamma^\nu).
\end{equation}
由 Clifford 关系 $\{\gamma^\mu,\gamma^\nu\}=2\eta^{\mu\nu}$，直接计算
\begin{align*}
[\gamma^\mu\gamma^\nu,\gamma^\rho]
&=\gamma^\mu\{\gamma^\nu,\gamma^\rho\}
-\{\gamma^\mu,\gamma^\rho\}\gamma^\nu\\
&=2\eta^{\nu\rho}\gamma^\mu-2\eta^{\mu\rho}\gamma^\nu,\\
[[\gamma^\mu,\gamma^\nu],\gamma^\rho]
&=4(\eta^{\nu\rho}\gamma^\mu-\eta^{\mu\rho}\gamma^\nu).
\end{align*}
因此 $\Sigma^{\mu\nu}=\frac{\ii}{4}[\gamma^\mu,\gamma^\nu]$ 满足所需交换子。
回代得到
\[
\frac{\ii}{2}\omega_{\mu\nu}[\Sigma^{\mu\nu},\gamma^\rho]
=\frac12\omega_{\mu\nu}
(\eta^{\mu\rho}\gamma^\nu-\eta^{\nu\rho}\gamma^\mu)
=\omega^\rho{}_\sigma\gamma^\sigma.
\]
这同时核对了协变式的矩阵顺序、符号与因子。
\end{proof}

\begin{proposition}[一阶场论的 Euler--Lagrange 方程]
若 $S=\int\mathcal L(\Phi,\partial\Phi)\dd^dx$ 对任意紧支场变分驻定，则每个分量满足\[\frac{\partial\mathcal L}{\partial\Phi^A}-\partial_\mu\frac{\partial\mathcal L}{\partial(\partial_\mu\Phi^A)}=0.\]
\end{proposition}

\begin{proof}
对作用量变分，分部积分把 $\delta\Phi$ 的导数移走，使变分只以 $\delta\Phi^A$（而非 $\partial_\mu\delta\Phi^A$）的形式出现。然后利用 $\delta\Phi^A$ 的任意性读出方程。

写出变分。场的任意变分 $\Phi^A\to\Phi^A+\delta\Phi^A$ 引起
\begin{equation*}
\delta S=\int \dd^dx\left[\frac{\partial\mathcal{L}}{\partial\Phi^A}\delta\Phi^A+\frac{\partial\mathcal{L}}{\partial(\partial_\mu\Phi^A)}\partial_\mu\delta\Phi^A\right].
\end{equation*}
第一项来自 $\mathcal{L}$ 对场的显式依赖，第二项来自 $\mathcal{L}$ 对场导数的依赖。

定义共轭动量。为简洁起见定义
\begin{equation*}
\Pi_A^\mu\equiv\frac{\partial\mathcal{L}}{\partial(\partial_\mu\Phi^A)}.
\end{equation*}
物理上 $\Pi_A^0$ 是场的正则动量，$\Pi_A^i$ 是流的空间分量。

分部积分。第二项中 $\partial_\mu\delta\Phi^A$ 是变分的导数。用 Leibniz 法则把它移到 $\Pi_A^\mu$ 上：
\begin{equation*}
\Pi_A^\mu\partial_\mu\delta\Phi^A=\partial_\mu(\Pi_A^\mu\delta\Phi^A)-(\partial_\mu\Pi_A^\mu)\delta\Phi^A.
\end{equation*}
第一项是全散度，积分后化为边界项。

代入并丢掉边界项。
\begin{equation*}
\delta S=\int \dd^dx\left[\left(\frac{\partial\mathcal{L}}{\partial\Phi^A}-\partial_\mu\Pi_A^\mu\right)\delta\Phi^A+\partial_\mu(\Pi_A^\mu\delta\Phi^A)\right].
\end{equation*}
若变分 $\delta\Phi^A$ 在边界上消失（或场在无穷远衰减足够快），散度项不贡献。

利用 $\delta\Phi^A$ 任意性。剩余积分中 $\delta\Phi^A$ 任意，要使 $\delta S=0$ 对所有 $\delta\Phi^A$ 成立，系数必须为零：
\begin{equation}\label{eq:ms-poincare-EL}
\boxed{
\frac{\partial\mathcal{L}}{\partial\Phi^A}-\partial_\mu\frac{\partial\mathcal{L}}{\partial(\partial_\mu\Phi^A)}=0.
}
\end{equation}
这就是 Euler--Lagrange 方程。对每个场分量 $A$ 有一个方程。

推广与边界：高阶导数 Lagrangian 给出高阶 Euler--Lagrange 方程，退化密度的约束需 Dirac--Bergmann 分析。
变分边界项 $\Theta^\mu=\Pi_A^\mu\delta\Phi^A$ 本身还不是一般 Noether 流。
若某对称变分满足 $\delta\mathcal L=\partial_\mu K^\mu$，则
$J^\mu=\Theta^\mu-K^\mu$ 在壳守恒，荷由 $Q=\int_\Sigma J^\mu\dd\Sigma_\mu$ 定义。
Euler--Lagrange 方程与 Noether 恒等式同源，但不能对任意变分把 $\Theta$ 直接当守恒流。
\end{proof}

\begin{proposition}[平移 Noether 流与正则能动量张量]
若一阶拉氏密度无显含坐标且 $\Pi_A^\mu=\partial\mathcal L/\partial(\partial_\mu\Phi^A)$，则\[T^\mu{}_\nu=\Pi_A^\mu\partial_\nu\Phi^A-\delta^\mu{}_\nu\mathcal L,\qquad\partial_\mu T^\mu{}_\nu=0\]在场方程成立时成立。
\end{proposition}

\begin{proof}
把平移变分 $\delta\Phi^A=\epsilon^\nu\partial_\nu\Phi^A$ 代入 Noether 流\eqnrefs{eq:ms-poincare-J-def}，提出任意参数 $\epsilon^\nu$，剩余部分定义为能动量张量。关键步骤是把 $K^\mu=\epsilon^\mu\mathcal{L}$ 中的指标 $\mu$ 改写为 $\epsilon^\nu$ 的缩并形式，使所有项共享同一个任意参数。

Symplectic 势。平移变分 $\delta\Phi^A=\epsilon^\nu\partial_\nu\Phi^A$ 代入 $\Theta^\mu=\Pi_A^\mu\delta\Phi^A$：
\begin{equation*}
\Theta^\mu=\Pi_A^\mu\,\epsilon^\nu\partial_\nu\Phi^A=\epsilon^\nu\Pi_A^\mu\partial_\nu\Phi^A.
\end{equation*}
这里 $\epsilon^\nu$ 是任意常矢量（平移参数），可以提到缩并外面。$\Pi_A^\mu\partial_\nu\Phi^A$ 是"动量密度乘速度"的结构——$\Pi_A^\mu$ 是共轭动量（场的 $\mu$ 方向流），$\partial_\nu\Phi^A$ 是场在 $\nu$ 方向的变化率。

Noether 流。由\eqnrefs{eq:ms-poincare-J-def}，$J^\mu=\Theta^\mu-K^\mu$。平移变换使拉氏量变分为全散度 $\delta\mathcal{L}=\epsilon^\nu\partial_\nu\mathcal{L}=\partial_\mu(\epsilon^\mu\mathcal{L})$（利用 $\epsilon^\nu$ 常矢量），故 $K^\mu=\epsilon^\mu\mathcal{L}$。Noether 流为
\begin{equation*}
J^\mu=\epsilon^\nu\Pi_A^\mu\partial_\nu\Phi^A-\epsilon^\mu\mathcal{L}.
\end{equation*}
第二项 $\epsilon^\mu\mathcal{L}$ 的自由指标是 $\mu$，而第一项的任意参数指标是 $\nu$。为了合并两项，需要统一任意参数。利用 $\epsilon^\mu=\delta^\mu{}_\nu\epsilon^\nu$（Kronecker $\delta$ 转指标）：
\begin{equation*}
\epsilon^\mu\mathcal{L}=\delta^\mu{}_\nu\epsilon^\nu\mathcal{L}=\epsilon^\nu\delta^\mu{}_\nu\mathcal{L}.
\end{equation*}
代入后 Noether 流统一为
\begin{equation*}
J^\mu=\epsilon^\nu\big[\Pi_A^\mu\partial_\nu\Phi^A-\delta^\mu{}_\nu\mathcal{L}\big].
\end{equation*}

定义 正则 能动量张量。因为 $\epsilon^\nu$ 任意，方括号中的量对每个 $\nu$ 独立守恒。定义 正则 能动量张量
\begin{equation}\label{eq:ms-poincare-T-can}
\boxed{
T^\mu{}_\nu=\Pi_A^\mu\partial_\nu\Phi^A-\delta^\mu{}_\nu\mathcal{L}.
}
\end{equation}
于是 $J^\mu_{\text{trans}}=\epsilon^\nu T^\mu{}_\nu$。第一项 $\Pi_A^\mu\partial_\nu\Phi^A$ 是"动量流乘场速度"——物理上是能量-动量流的密度。第二项 $-\delta^\mu{}_\nu\mathcal{L}$ 是拉氏密度的对角项——它来自平移变分使拉氏量本身移动的效应。

守恒律。由 Noether 定理\eqnrefs{eq:ms-poincare-noether}，$\partial_\mu J^\mu=0$ 在壳。代入 $J^\mu=\epsilon^\nu T^\mu{}_\nu$：
\begin{equation*}
0=\partial_\mu J^\mu=\partial_\mu(\epsilon^\nu T^\mu{}_\nu)=\epsilon^\nu\,\partial_\mu T^\mu{}_\nu.
\end{equation*}
最后一步利用 $\epsilon^\nu$ 常矢量（$\partial_\mu\epsilon^\nu=0$）。由于 $\epsilon^\nu$ 任意，每个分量独立为零：
\begin{equation}\label{eq:ms-poincare-T-conservation}
\boxed{\partial_\mu T^\mu{}_\nu=0.}
\end{equation}

一般 $T^{\mu\nu}$ 不对称，但可加全导数项（Belinfante 修正）使 $\Theta^{\mu\nu}$ 对称且守恒，详见下框。
能动量张量也可由作用量对度量的变分定义 $T_{\mu\nu}=\frac{-2}{\sqrt{-g}}\frac{\delta S}{\delta g^{\mu\nu}}$，自动给出对称张量，并在弯曲时空推广为协变守恒 $\nabla_\mu T^{\mu\nu}=0$。
\end{proof}

\begin{proposition}[Lorentz 总角动量流守恒]
对 Lorentz 不变的一阶作用量，按本章约定\[L^{\lambda\mu\nu}=x^\mu T^{\lambda\nu}-x^\nu T^{\lambda\mu},\qquad S^{\lambda\mu\nu}=-\ii\Pi_A^\lambda(\Sigma^{\mu\nu})^A{}_B\Phi^B.\]总流满足 $\partial_\lambda(L+S)^{\lambda\mu\nu}=0$，且 $\partial_\lambda S^{\lambda\mu\nu}=T^{\nu\mu}-T^{\mu\nu}$。
\end{proposition}

\begin{proof}
把 Lorentz 变分代入 Noether 流，分离出轨道部分（含 $x$ 和 $T$）和自旋部分（含 $\Sigma$），分别定义轨道角动量流和自旋流。

Noether 流。
\begin{equation*}
J^\lambda=\Pi_A^\lambda\delta\Phi^A-\xi^\lambda\mathcal{L}
=\omega^\rho{}_\sigma x^\sigma\big(\Pi_A^\lambda\partial_\rho\Phi^A-\delta^\lambda{}_\rho\mathcal{L}\big)+\frac{i}{2}\omega_{\mu\nu}\Pi_A^\lambda(\Sigma^{\mu\nu})^A{}_B\Phi^B.
\end{equation*}

识别能动量张量。第一括号正是 $T^\lambda{}_\rho$（见\eqnrefs{eq:ms-poincare-T-can}）。

轨道部分改写。利用 $\omega_{\mu\nu}$ 反对称性，把 $\omega^\rho{}_\sigma x^\sigma T^\lambda{}_\rho$ 改写为 $\omega_{\mu\nu}$ 的缩并：
\begin{equation*}
\omega^\rho{}_\sigma x^\sigma T^\lambda{}_\rho=-\frac{1}{2}\omega_{\mu\nu}(x^\mu T^{\lambda\nu}-x^\nu T^{\lambda\mu}).
\end{equation*}
这与\eqnrefs{eq:ms-poincare-L-def} 中轨道部分的改写完全平行。

定义轨道与自旋流。
\begin{equation}\label{eq:ms-poincare-L-current}
L^{\lambda\mu\nu}=x^\mu T^{\lambda\nu}-x^\nu T^{\lambda\mu},
\end{equation}
\begin{equation}\label{eq:ms-poincare-S-current}
S^{\lambda\mu\nu}=-i\,\Pi_A^\lambda(\Sigma^{\mu\nu})^A{}_B\Phi^B.
\end{equation}
轨道流 $L^{\lambda\mu\nu}$ 只依赖坐标和能动量张量，对所有场类型一样。自旋流 $S^{\lambda\mu\nu}$ 依赖场的内禀表示 $\Sigma^{\mu\nu}$，标量场为零。

合并。
\begin{equation*}
J^\lambda_{\text{Lorentz}}=-\frac{1}{2}\omega_{\mu\nu}(L^{\lambda\mu\nu}+S^{\lambda\mu\nu}).
\end{equation*}
定义总角动量流
\begin{equation}\label{eq:ms-poincare-J-total}
J^{\lambda\mu\nu}=L^{\lambda\mu\nu}+S^{\lambda\mu\nu},
\end{equation}
由 $\omega_{\mu\nu}$ 任意得
\begin{equation}\label{eq:ms-poincare-J-conservation}
\boxed{\partial_\lambda J^{\lambda\mu\nu}=0.}
\end{equation}

对上面的轨道流求散度，并使用 $\partial_\lambda T^{\lambda\nu}=0$，得
\[
\partial_\lambda L^{\lambda\mu\nu}=T^{\mu\nu}-T^{\nu\mu},\qquad
\partial_\lambda S^{\lambda\mu\nu}=T^{\nu\mu}-T^{\mu\nu}.
\]
所以正则轨道流和自旋流通常只有总和守恒。
按本章变分约定，自旋流始终为
$S^{\lambda\mu\nu}=-\ii\Pi_A^\lambda(\Sigma^{\mu\nu})^A{}_B\Phi^B$。
标量场的生成元为零，故自旋流为零。
对于采用一阶 Dirac 密度的旋量场，$\Pi_\psi^\lambda=\ii\bar\psi\gamma^\lambda$，
代入得到 $S^{\lambda\mu\nu}=\bar\psi\gamma^\lambda\Sigma^{\mu\nu}\psi$。
若改用对称化密度，正则流也必须相应重算，不能混用两种表达。
矢量场同样应从其共轭动量和矢量生成元代入，而非猜测自旋项的符号。

Belinfante 改进同时重分配轨道和自旋流，使总角动量可用改进后的对称张量表达；
它不意味着原来那个正则自旋流自动单独守恒。
量子化后总角动量满足旋转代数；$P^2$ 与 Pauli--Lubanski 矢量平方 $W^2$
再结合能量分支和小群表示用于 Wigner 分类。
对无质量表示，$P^2=W^2=0$ 不能区分所有螺旋度，不能宣称两个 Casimir
单独完全标记全部单粒子表示。
\end{proof}

\begin{proposition}[Belinfante 张量的对称性与守恒性]
设 $T,S$ 满足前述 Noether 恒等式。定义\[B^{\lambda\mu\nu}=\tfrac12(S^{\lambda\mu\nu}+S^{\mu\nu\lambda}+S^{\nu\mu\lambda}),\qquad T_B^{\mu\nu}=T^{\mu\nu}+\partial_\lambda B^{\lambda\mu\nu}.\]则 $T_B$ 对称且守恒；改进前后总角动量流只差全散度，适当边界衰减下总荷相同。
\end{proposition}

\begin{proof}
利用反对称关系\eqnrefs{eq:ms-poincare-T-asym} 与 Belinfante 修正\eqnrefs{eq:ms-poincare-B-def} 的指标性质，分三步验证：(1) $T_B^{\mu\nu}=T_B^{\nu\mu}$（对称性），(2) $\partial_\mu T_B^{\mu\nu}=0$（守恒性），(3) 角动量流可用 $T_B$ 单独表示（自旋吸收）。

先计算 $T_B^{\mu\nu}-T_B^{\nu\mu}$。由定义\eqnrefs{eq:ms-poincare-TB-def}，
\begin{equation*}
T_B^{\mu\nu}-T_B^{\nu\mu}=(T_{\text{can}}^{\mu\nu}-T_{\text{can}}^{\nu\mu})+\partial_\lambda(B^{\lambda\mu\nu}-B^{\lambda\nu\mu}).
\end{equation*}
第一项由\eqnrefs{eq:ms-poincare-T-asym} 给出：$T_{\text{can}}^{\mu\nu}-T_{\text{can}}^{\nu\mu}=-\partial_\lambda S^{\lambda\mu\nu}$。第二项需计算 $B^{\lambda\mu\nu}-B^{\lambda\nu\mu}$。由\eqnrefs{eq:ms-poincare-B-def}：
\begin{equation*}
B^{\lambda\mu\nu}=\frac{1}{2}\big(S^{\lambda\mu\nu}+S^{\mu\nu\lambda}+S^{\nu\mu\lambda}\big),\qquad
B^{\lambda\nu\mu}=\frac{1}{2}\big(S^{\lambda\nu\mu}+S^{\nu\mu\lambda}+S^{\mu\nu\lambda}\big).
\end{equation*}
相减时 $S^{\mu\nu\lambda}$ 和 $S^{\nu\mu\lambda}$ 各自抵消：
\begin{equation*}
B^{\lambda\mu\nu}-B^{\lambda\nu\mu}=\frac{1}{2}\big(S^{\lambda\mu\nu}-S^{\lambda\nu\mu}\big).
\end{equation*}
利用自旋流后两指标反对称性 $S^{\lambda\nu\mu}=-S^{\lambda\mu\nu}$：
\begin{equation*}
B^{\lambda\mu\nu}-B^{\lambda\nu\mu}=\frac{1}{2}\big(S^{\lambda\mu\nu}+S^{\lambda\mu\nu}\big)=S^{\lambda\mu\nu}.
\end{equation*}
因此
\begin{equation*}
T_B^{\mu\nu}-T_B^{\nu\mu}=-\partial_\lambda S^{\lambda\mu\nu}+\partial_\lambda S^{\lambda\mu\nu}=0.
\end{equation*}
Belinfante 张量对称。

守恒性。由\eqnrefs{eq:ms-poincare-TB-def}，$\partial_\mu T_B^{\mu\nu}=\partial_\mu T_{\text{can}}^{\mu\nu}+\partial_\mu\partial_\lambda B^{\lambda\mu\nu}$。第一项为零（正则 守恒\eqnrefs{eq:ms-poincare-T-conservation}）。第二项中 $B^{\lambda\mu\nu}$ 对前两指标反对称：$B^{\lambda\mu\nu}=-B^{\mu\lambda\nu}$。验证：由定义
\begin{equation*}
B^{\mu\lambda\nu}=\frac{1}{2}\big(S^{\mu\lambda\nu}+S^{\lambda\nu\mu}+S^{\nu\lambda\mu}\big),
\end{equation*}
只使用后两指标的反对称性：$S^{\mu\lambda\nu}=-S^{\mu\nu\lambda}$、
$S^{\lambda\nu\mu}=-S^{\lambda\mu\nu}$、
$S^{\nu\lambda\mu}=-S^{\nu\mu\lambda}$，重新排列三项即可验证 $B^{\mu\lambda\nu}=-B^{\lambda\mu\nu}$。因此 $\partial_\mu\partial_\lambda$（对称导数）乘 $B^{\lambda\mu\nu}$（前两指标反对称）恒为零：
\begin{equation*}
\partial_\mu\partial_\lambda B^{\lambda\mu\nu}=\frac{1}{2}(\partial_\mu\partial_\lambda-\partial_\lambda\partial_\mu)B^{\lambda\mu\nu}=0.
\end{equation*}
故 $\partial_\mu T_B^{\mu\nu}=0+0=0$。

角动量流改写。用 $T_B$ 改写总"纯轨道"角动量流：
\begin{equation*}
J_B^{\lambda\mu\nu}=x^\mu T_B^{\lambda\nu}-x^\nu T_B^{\lambda\mu}.
\end{equation*}
改进前后的差可逐项写成
\[
J_B^{\lambda\mu\nu}-(L^{\lambda\mu\nu}+S^{\lambda\mu\nu})
=\partial_\rho\bigl(x^\mu B^{\rho\lambda\nu}-x^\nu B^{\rho\lambda\mu}\bigr),
\]
因为 $B^{\mu\lambda\nu}-B^{\nu\lambda\mu}=-S^{\lambda\mu\nu}$。
故在空间无穷远衰减足够快或边界通量消失时，两者总荷相同。
这个形式的角动量流守恒 $\partial_\lambda J_B^{\lambda\mu\nu}=0$ 等价于原来的 $\partial_\lambda(L^{\lambda\mu\nu}+S^{\lambda\mu\nu})=0$，但自旋信息已被吸收到对称的 $T_B$ 中——不再需要单独的自旋流。这就是 Belinfante 改进的物理意义：把"轨道+自旋"的分解重新打包成一个对称张量的"纯轨道"形式。
\end{proof}
